\section{Related Work}
\label{sec:related-work}

Classical cloud simulators, such as CloudSim~\cite{calheiros2011cloudsim}, iCanCloud~\cite{nunez2012icancloud}, and SimGrid~\cite{casanova2008simgrid}, are designed to evaluate resource management and scheduling without physical infrastructure~\cite{jawaddi2024gymcloudsim}. Many of these foundational tools have been extended to support specific operational needs; for example, WorkflowSim~\cite{chen2012workflowsim} adds complex scientific workflow modeling, while frameworks like GreenCloud~\cite{kliazovich2012greencloud} and E-mc$^2$~\cite{castane2013emc2} introduce formal energy modeling for data centers~\cite{makaratzis2018energymodeling}. Collectively, they demonstrate the utility of discrete event simulation for analyzing system dynamics and computational accounting.

As quantum computing increasingly adopts a cloud-based delivery model dedicated simulation toolkits have arisen to model hybrid quantum environments. Frameworks such as iQuantum~\cite{nguyen2024iquantum}, QSimPy~\cite{qsimpy2024, nguyen2024drlq}, and QCloudSim~\cite{luo2025digital} apply discrete event simulation to evaluate quantum task scheduling, backend selection, and resource management strategies. 

While simulation frameworks are critical for exploring these emerging architectures, existing quantum cloud simulators~\cite{Jones2019-nc, Brennan2022QXToolsAJ, diadamo2021qunetsim} focus primarily on performance, resource allocation, and queueing. They generally lack explicit energy and cost estimation at the workflow level. Furthermore, many frameworks~\cite{bian2023pas, fan2021drascqsim} treat quantum jobs as monolithic, overlooking the interleaved QPU and CPU execution phases central to hybrid algorithms.

Prior work shows that practical quantum overhead is driven by workload traits, backend variability, repeated sampling, and classical-quantum interactions. Senapati et al.~\cite{senapati2023reproducibility} showed hardware calibration drift degrades workload reproducibility across backends. Related research characterized noise-sensitive states~\cite{senapati2023visualization,senapati2025noisy} and devised predictive NISQ backend selection for quantum machine learning~\cite{senapati2024pqml}. Subsequent uncertainty quantification studies analyzed execution variability, stability, and backend selection for variational quantum algorithms and VQE~\cite{senapati2025uqvarqa,senapati2026vqe}. Additional work across hybrid molecular simulation~\cite{abraham2026manybody}, machine learning translation~\cite{zeng2026qbridge}, and variational/signal processing workloads~\cite{senapati2026unified} confirms that total execution cost is determined by the end-to-end hybrid workflow---including iterative classical optimization, repeated executions, and backend choice---rather than isolated quantum circuits.

Energy consumption is a critical concern in high performance computing (HPC)~\cite{Veigas2025, iea2025data_centers}, driving interest in quantum computing as a potentially energy efficient complement~\cite{martin2022energy, jaschke2023quantum}. While classical supercomputers can consume hundreds of megawatt-hours daily, present day quantum processors operate in the kilowatt range. However, meaningful energy evaluation requires analyzing total system power integrated over execution time, as a quantum computer's footprint is dictated by its supporting infrastructure rather than gate level operations. As summarized in Figure~\ref{fig:energy_usage}~\cite{ezratty2021}, superconducting systems demand significantly more power due to their environmental control requirements than other quantum systems.

In superconducting platforms, energy consumption is overwhelmingly dominated by system level infrastructure. Total system power ranges from 25~kW to over 140~kW, depending on system configuration and qubit count~\cite{ezratty2021}. The primary contributor is the cryogenic refrigeration required to maintain millikelvin temperatures, drawing between 16~kW and 105~kW. Additional components, such as classical computing overhead and qubit gate control electronics, contribute marginally (often around 1~kW) and are generally subsumed into the total system consumption.

Based on this observation, we developed an energy accounting model in which the energy consumed by a quantum job is proportional to the duration for which it occupies QPUs and CPUs. Under this abstraction, the energy consumption of a superconducting quantum job is approximated as $E_{\text{job}} \approx P_{\text{baseline}} \cdot t_{\text{occupy}}$, where \(P_{\text{baseline}}\) denotes the average system-level power consumption, including cryogenics and supporting electronics, and \(t_{\text{occupy}}\) represents the job's execution interval on the quantum device. This baseline power model aligns with energy accounting practices commonly used in HPC systems, where always-on infrastructure costs are amortized over job execution time. 